\begin{abstract}
We investigate why biological attention mechanisms benefit some neural architectures more than others. We introduce the \textit{compression-covariance dichotomy}: RNNs and LSTMs compress sequence history into fixed-size states, creating an information bottleneck, while transformers compute attention as an explicit covariance-like operation that preserves all pairwise relationships. We implement four neuroscience-inspired mechanisms---biased competition (top-down task bias), lateral inhibition (competitive suppression), tuning sharpening (enhanced selectivity), and noise decorrelation (multi-head diversity)---across RNN, LSTM, and Transformer architectures. Our experiments reveal a clear hierarchy: transformers improve by 11.7\%, LSTM by 6.2\%, and RNN by 4.2\%, directly reflecting each architecture's capacity to preserve relational structure. Bio-enhanced transformers exhibit faster convergence (15--20\% speedup), emergent attention sparsity (23\% $\rightarrow$ 41\% sparse weights), reduced inter-head correlation (0.38 $\rightarrow$ 0.18), and more uniform gradient flow across layers. Ablation studies show that lateral inhibition and tuning sharpening contribute most, with mechanisms acting complementarily. These findings establish that biological attention mechanisms are most effective when applied to architectures that preserve, rather than compress, relational structure---a principle with implications for future architecture design.
\end{abstract}
